\section{Snap-HyRE}\label{sec:method}

\subsection{Design Rationale}

\hyre{} mirrors a small piece of how an expert reads a legal scenario. A
practitioner facing a fact pattern forms a quick intuition about what is at
issue (proximate cause? a hearsay exception? a state-specific notice rule?),
looks up the controlling authority by that intuition, and then judges the
evidence on its own merits rather than anchoring on the initial guess.
\hyre{} encodes the same loop. The first call is the intuition step: it
produces a draft answer and a separate legal search passage written in the
style of the controlling authority. The second call is the evidence step:
it sees only the original question and the retrieved evidence, never the
draft, so the final answer is judged against what the corpus actually says.
Figure~\ref{fig:pipeline_art} shows the resulting two-call pipeline.

\subsection{Two Calls, One Retrieval Query}\label{sec:method:snaphyre}

\hyre{} separates query formation from final answering. The same model
$G_\theta$ is used for both calls; only the prompt template differs (the
superscript). Given a formatted legal question $q$, the first call produces
a private draft answer $a_0$ and a legal search passage $p$:

\begin{equation}
(a_0, p) = G_{\theta}^{\mathrm{snap}}(q).
\end{equation}

The retriever $R_k$ embeds only $p$ and returns the top-$k$ retrieved
passages $E_k = R_k(p)$. The final answer is then

\begin{equation}
\hat{a} = G_{\theta}^{\mathrm{answer}}(q, E_k),
\end{equation}

with $a_0$ withheld from the final answer context. The first call can use
the model's prior knowledge of law to choose vocabulary; the second call
must answer from the retrieved corpus text.

Each \hyre{} row follows three steps: validate that the first call produced
both a formatted draft answer and a clean search passage; retrieve and rerank
using only the search passage; answer the original question from the retrieved
passages without exposing the draft answer. The retrieval passage is
constrained to be short and reference-like: it should describe the controlling
rule, statutory condition, legal relationship, or doctrinal test without
answer labels or an \texttt{Answer:} header.

The difference from HyDE is narrow but consequential. HyDE writes a legal
passage directly from the question; \hyre{} first forms a tentative legal
answer and uses that private judgment only to decide what kind of passage to
search with. For a torts fact pattern about a broken leg, crutches, and a
later fall, this turns an accident-heavy query into a passage about proximate
cause and foreseeable subsequent injury
(Figure~\ref{fig:query_contrast_diagram}).

\begin{figure*}[t]
  \centering
  \includegraphics[width=\textwidth]{22_query_contrast_diagram.png}
  \caption{What the two methods send to the retriever for one BarExamQA item
  (\texttt{mbe\_1175}, Gemma 4 26B). Raw-question RAG embeds the narrative
  fact pattern; \hyre{} embeds a short doctrinal passage about foreseeable
  intervening injury, which is the language of the rule the retriever
  actually needs to surface.}
  \label{fig:query_contrast_diagram}
\end{figure*}

\subsection{Why It Works: Two Guardrails}

The two-call structure gives the method two guardrails that follow from the
formulation above. First, the draft $a_0$ is private: it can shape what to
look for, but it is not handed in as evidence, so the second call avoids the
anchoring bias of treating its own first guess as a fact. Second, the search
passage $p$ is written like legal authority rather than chain-of-thought
reasoning: it names the rule or statutory condition that would resolve the
question, not the steps the model used to get there. The intuition step is
allowed to use model priors; the evidence step is forced back to the corpus.
